\documentclass[10pt]{article}
\newcommand{\DocShort}{Rúbrica del examen de hidrostática}
% Layout specific to the one-page landscape grading rubric.
\input{exam_preamble.tex}

\geometry{letterpaper, landscape, top=1.45cm, bottom=1.45cm, left=1.45cm, right=1.45cm}

\fancypagestyle{rubric}{
  \fancyhf{}
  \fancyhead[L]{\small\color{midgrey}ICV-343-T Hidráulica Básica}
  \fancyhead[R]{\small\color{midgrey}\DocShort}
  \fancyfoot[C]{\small\color{midgrey}\thepage}
  \renewcommand{\headrulewidth}{0.4pt}
}
\pagestyle{rubric}

% The shared preamble is portrait by default. Recalculate fancyhdr after the
% landscape geometry is active so the header rule spans the full text width.
\AtBeginDocument{\setlength{\headwidth}{\textwidth}}

\setlength{\parindent}{0pt}
\newcommand{\nivelcasilla}{\centering\Large$\square$}

\begin{document}
\begin{center}
{\Large\bfseries\color{darkblue}Rúbrica de corrección}\par
{\normalsize\bfseries\color{darkblue}Examen de hidrostática y propiedades seleccionadas de los fluidos}\par
\vspace{0.06cm}
{\color{medblue}\rule{\linewidth}{1.3pt}}
\end{center}

\vspace{0.05cm}
\begin{tabularx}{\linewidth}{@{}>{\bfseries}l X >{\bfseries}l p{4.5cm}@{}}
Estudiante: & \hrulefill & Matrícula: & \hrulefill \\
Profesor: & Ricardo Rafael Hernández Moreira, PhD & Fecha: & \hrulefill
\end{tabularx}

\vspace{0.12cm}
\renewcommand{\arraystretch}{1.16}
\rowcolors{2}{white}{white}
\small
\begin{tabularx}{\linewidth}{
  >{\centering\arraybackslash}p{1.75cm}
  >{\raggedright\arraybackslash}X
  >{\raggedright\arraybackslash}X
  >{\raggedright\arraybackslash}X
  >{\raggedright\arraybackslash}X
  >{\centering\arraybackslash}p{2.25cm}}
\toprule
\rowcolor{lightblue}
\textbf{Problema} &
\textbf{Dominio sobresaliente}\newline
Solución completa, clara y trazable: esquema, datos, principios físicos, desarrollo en SI, resultado e interpretación correctos.\newline
\textbf{90--100\%} &
\textbf{Dominio competente}\newline
Método esencialmente correcto y bien sustentado; presenta omisiones o errores menores que no alteran el enfoque principal.\newline
\textbf{75--89\%} &
\textbf{Dominio en desarrollo}\newline
Reconoce parte del modelo y avanza con un planteamiento pertinente, pero incurre en errores u omisiones importantes.\newline
\textbf{50--74\%} &
\textbf{Dominio inicial o insuficiente}\newline
Evidencia limitada, enfoque inapropiado o procedimiento no trazable; resultados sin sustento suficiente.\newline
\textbf{0--49\%} &
\textbf{Calificación} \\
\midrule
1 {\normalfont\footnotesize (15 pt)} & \nivelcasilla & \nivelcasilla & \nivelcasilla & \nivelcasilla & \rule{0pt}{0.88cm}\textbf{\hspace{0.7cm}/ 15} \\
\midrule
2 {\normalfont\footnotesize (20 pt)} & \nivelcasilla & \nivelcasilla & \nivelcasilla & \nivelcasilla & \rule{0pt}{0.88cm}\textbf{\hspace{0.7cm}/ 20} \\
\midrule
3 {\normalfont\footnotesize (15 pt)} & \nivelcasilla & \nivelcasilla & \nivelcasilla & \nivelcasilla & \rule{0pt}{0.88cm}\textbf{\hspace{0.7cm}/ 15} \\
\midrule
4 {\normalfont\footnotesize (25 pt)} & \nivelcasilla & \nivelcasilla & \nivelcasilla & \nivelcasilla & \rule{0pt}{0.88cm}\textbf{\hspace{0.7cm}/ 25} \\
\midrule
5 {\normalfont\footnotesize (12.5 pt)} & \nivelcasilla & \nivelcasilla & \nivelcasilla & \nivelcasilla & \rule{0pt}{0.88cm}\textbf{\hspace{0.45cm}/ 12.5} \\
\midrule
6 {\normalfont\footnotesize (12.5 pt)} & \nivelcasilla & \nivelcasilla & \nivelcasilla & \nivelcasilla & \rule{0pt}{0.88cm}\textbf{\hspace{0.45cm}/ 12.5} \\
\bottomrule
\end{tabularx}

\vspace{0.12cm}
\begin{tabularx}{\linewidth}{@{}>{\bfseries}p{3.2cm} X >{\bfseries\centering\arraybackslash}p{4.2cm}@{}}
Observaciones generales: & \rule{0pt}{0.55cm}\hrulefill & Calificación final: \rule{1.2cm}{0.4pt}\,/\,100
\end{tabularx}

\end{document}
